% !TEX program = xelatex
% Lernplatt - by Can Siebert
% Kurs 71 (Dig 42) - Tag 10 - Loesungsheft
% Skalierbare digitale Prozessarchitektur: Microservices, Kubernetes, TDD/DevOps, Operating Model
%
% Self-contained xelatex source. Kein biblatex, keine externen Bilder.

\documentclass[11pt,a4paper,oneside]{article}

% --- Encoding & Sprache --------------------------------------------------
\usepackage{polyglossia}
\setdefaultlanguage{german}

% --- Geometrie -----------------------------------------------------------
\usepackage[a4paper,margin=2.5cm]{geometry}

% --- Fonts ---------------------------------------------------------------
\usepackage{fontspec}
\defaultfontfeatures{Ligatures=TeX}

\IfFontExistsTF{Charter}{%
  \setmainfont{Charter}%
}{%
  \IfFontExistsTF{Noto Serif}{%
    \setmainfont{Noto Serif}%
  }{%
    \setmainfont{Liberation Serif}%
  }%
}

\IfFontExistsTF{Source Sans 3}{%
  \setsansfont{Source Sans 3}%
}{%
  \IfFontExistsTF{Source Sans Pro}{%
    \setsansfont{Source Sans Pro}%
  }{%
    \setsansfont{Liberation Sans}%
  }%
}

\newcommand{\caveatfontfallback}{\itshape}
\IfFontExistsTF{Caveat}{%
  \newfontfamily\caveatfont{Caveat}[Scale=1.15]%
}{%
  \newcommand\caveatfont{\caveatfontfallback}%
}

\setmonofont{Liberation Mono}

% --- Mikro-Typografie ----------------------------------------------------
\usepackage{microtype}
\usepackage{eurosym}
\linespread{1.15}

% --- Farben & Boxen ------------------------------------------------------
\usepackage{xcolor}
\definecolor{lpInk}{RGB}{20,28,38}
\definecolor{lpAccent}{RGB}{157,74,26}
\definecolor{lpAccentLight}{RGB}{246,228,212}
\definecolor{lpAmber}{RGB}{222,138,30}
\definecolor{lpAmberLight}{RGB}{255,243,217}
\definecolor{lpAmberBorder}{RGB}{196,118,18}
\definecolor{lpGray}{RGB}{96,104,114}
\definecolor{lpGrayLight}{RGB}{238,240,243}
\definecolor{lpGreen}{RGB}{34,120,72}
\definecolor{lpGreenLight}{RGB}{222,238,228}
\definecolor{lpGreenBorder}{RGB}{24,96,58}

\definecolor{lpL1}{RGB}{34,120,72}
\definecolor{lpL2}{RGB}{22,90,140}
\definecolor{lpL3}{RGB}{170,42,52}

\usepackage[most]{tcolorbox}
\tcbuselibrary{breakable,skins}

\newtcolorbox{infobox}[1][Hinweis]{%
  enhanced, breakable,
  colback=lpAccentLight,
  colframe=lpAccent,
  boxrule=0.6pt,
  arc=2pt,
  left=10pt,right=10pt,top=8pt,bottom=8pt,
  fonttitle=\sffamily\bfseries\color{white},
  coltitle=white,
  title={#1},
  attach boxed title to top left={xshift=8pt,yshift=-8pt},
  boxed title style={size=small, colback=lpAccent, colframe=lpAccent, boxrule=0.4pt, arc=1pt},
  top=14pt
}

\newtcolorbox{loesungbox}[1][Musterl\"osung]{%
  enhanced, breakable,
  colback=lpGreenLight,
  colframe=lpGreenBorder,
  boxrule=0.6pt,
  arc=2pt,
  left=10pt,right=10pt,top=8pt,bottom=8pt,
  fonttitle=\sffamily\bfseries\color{white},
  coltitle=white,
  title={#1},
  attach boxed title to top left={xshift=8pt,yshift=-8pt},
  boxed title style={size=small, colback=lpGreenBorder, colframe=lpGreenBorder, boxrule=0.4pt, arc=1pt},
  top=14pt
}

% --- Listen --------------------------------------------------------------
\usepackage{enumitem}
\setlist{itemsep=2pt, topsep=4pt, parsep=0pt}
\setlist[itemize,1]{label={\textbullet},leftmargin=1.6em}
\setlist[itemize,2]{label={\textendash},leftmargin=1.4em}
\setlist[enumerate,1]{label=\arabic*., leftmargin=1.8em}

% --- Tabellen ------------------------------------------------------------
\usepackage{tabularx}
\usepackage{array}
\usepackage{booktabs}
\usepackage{longtable}

% --- Kopf-/Fusszeilen ----------------------------------------------------
\usepackage{fancyhdr}
\usepackage{lastpage}
\pagestyle{fancy}
\fancyhf{}
\renewcommand{\headrulewidth}{0.3pt}
\renewcommand{\footrulewidth}{0pt}
\fancyhead[L]{\footnotesize\sffamily\color{lpGray}L\"osungsheft \textperiodcentered{} Kurs 71 \textperiodcentered{} Tag 10}
\fancyhead[R]{\footnotesize\sffamily\color{lpGray}Microservices, Kubernetes \& DevOps}
\fancyfoot[L]{\footnotesize\sffamily\color{lpGray}\textcopyright{} Can Siebert}
\fancyfoot[R]{\footnotesize\sffamily\color{lpGray}\thepage{} / \pageref{LastPage}}

\fancypagestyle{plain}{%
  \fancyhf{}%
  \renewcommand{\headrulewidth}{0pt}%
  \fancyfoot[C]{\footnotesize\sffamily\color{lpGray}\thepage{} / \pageref{LastPage}}%
}

% --- Section-Styling -----------------------------------------------------
\usepackage[explicit]{titlesec}

\titleformat{\section}[block]
  {\sffamily\Large\bfseries\color{lpAccent}}
  {\thesection.}{0.6em}{#1}
  [{\vspace{2pt}\color{lpAccent}\titlerule[0.6pt]}]

\titleformat{\subsection}[block]
  {\sffamily\large\bfseries\color{lpInk}}
  {\thesubsection}{0.6em}{#1}

\titlespacing*{\section}{0pt}{14pt}{8pt}
\titlespacing*{\subsection}{0pt}{10pt}{4pt}

% --- Hyperlinks ----------------------------------------------------------
\usepackage[hidelinks,unicode]{hyperref}

% --- Helfer --------------------------------------------------------------
\newcommand{\akzent}[1]{{\caveatfont\color{lpAmber}#1}}
\newcommand{\fachbegriff}[1]{\textbf{#1}}

\newcommand{\niveaubadge}[2]{%
  \tcbox[on line, colback=#1, colframe=#1, coltext=white,
         boxrule=0pt, arc=2pt, left=5pt,right=5pt,top=1pt,bottom=1pt,
         nobeforeafter]{\sffamily\bfseries\footnotesize #2}%
}
\newcommand{\nivLi}{\niveaubadge{lpL1}{L1\,\textbullet\,Wissen}}
\newcommand{\nivLii}{\niveaubadge{lpL2}{L2\,\textbullet\,Anwendung}}
\newcommand{\nivLiii}{\niveaubadge{lpL3}{L3\,\textbullet\,Management}}

\newcommand{\zeit}[1]{%
  \tcbox[on line, colback=lpGrayLight, colframe=lpGrayLight, coltext=lpInk,
         boxrule=0pt, arc=2pt, left=5pt,right=5pt,top=1pt,bottom=1pt,
         nobeforeafter]{\sffamily\footnotesize\textbf{$\sim$\,#1}}%
}

\newcommand{\aufgabenkopf}[4]{%
  \par\vspace{6pt}
  \noindent{\sffamily\Large\bfseries\color{lpAccent}Aufgabe #1 \textendash{} #2}\par
  \vspace{2pt}
  \noindent{\color{lpAccent}\rule{\linewidth}{0.6pt}}\par
  \vspace{4pt}
  \noindent #3 \hfill \zeit{#4}\par
  \vspace{6pt}
}

\newcommand{\ytquery}[1]{%
  \par\vspace{4pt}
  \noindent{\footnotesize\sffamily\color{lpGray}\textbf{YouTube-Suchquery:} \texttt{#1}}\par
}

% =========================================================================
\begin{document}

% --- COVER ---------------------------------------------------------------
\begin{titlepage}
  \pagestyle{empty}
  \vspace*{1.5cm}
  \noindent{\sffamily\footnotesize\color{lpGray}LERNPLATT \textperiodcentered{} BY CAN SIEBERT}\\[2pt]
  \noindent{\color{lpAccent}\rule{\linewidth}{1.2pt}}\\[4pt]
  \noindent{\sffamily\footnotesize\color{lpGray}L\"OSUNGSHEFT \textperiodcentered{} MODUL 4 \textperiodcentered{} TAG 10 VON 16 \textperiodcentered{} KURS 71 (Dig 42) \textperiodcentered{} 11.\,Juni 2026}

  \vspace{3cm}
  {\sffamily\fontsize{30}{34}\selectfont\bfseries\color{lpInk}L\"osungsheft\\Tag 10\par}

  \vspace{0.7cm}
  {\sffamily\large\color{lpGray}Musterl\"osungen zu den elf Aufgaben\\des Aufgabenhefts \textendash{} Microservices,\\Kubernetes, TDD und Operating Model.\par}

  \vspace{1.5cm}
  {\caveatfont\LARGE\color{lpGreen}Musterl\"osungen \textperiodcentered{} nicht die einzige Antwort}\par

  \vfill

  \noindent\begin{minipage}{0.55\linewidth}
    {\sffamily\footnotesize\color{lpGray}Hinweis:}\\
    {\sffamily\small\color{lpInk}Musterl\"osungen sind Diskussionsbasis,\\nicht die einzige korrekte L\"osung.}\\[10pt]
    {\sffamily\footnotesize\color{lpGray}Begleitend:}\\
    {\sffamily\small\color{lpInk}Aufgabenheft \& Lehrskript Tag 10}
  \end{minipage}\hfill
  \begin{minipage}{0.4\linewidth}
    \raggedleft
    {\sffamily\footnotesize\color{lpGray}Dozent:}\\
    {\sffamily\small\color{lpInk}Can Siebert}\\[10pt]
    {\sffamily\footnotesize\color{lpGray}Niveau-Mix:}\\
    {\sffamily\small\color{lpInk}3\,L1 \textperiodcentered{} 5\,L2 \textperiodcentered{} 3\,L3}
  \end{minipage}

  \vspace{0.5cm}
  \noindent{\color{lpAccent}\rule{\linewidth}{0.6pt}}
\end{titlepage}

% --- Anleitung -----------------------------------------------------------
\section*{Hinweise zu den Musterl\"osungen}
\thispagestyle{plain}

\noindent Die folgenden Musterl\"osungen sind \emph{eine} m\"ogliche, didaktisch nachvollziehbare L\"osung \textendash{} sie sind nicht die einzig richtige. Auf L2 und L3 sind mehrere Begr\"undungswege legitim, solange sie:

\begin{itemize}
  \item den im Lehrskript eingef\"uhrten Begriffsapparat (Microservices, Kubernetes, TDD, DevOps, Operating Model) sauber nutzen,
  \item Annahmen explizit machen, statt sie zu verschleiern,
  \item Zielkonflikte (Speed vs.\ Stabilit\"at, Autonomie vs.\ Standardisierung, Aufwand vs.\ Auditierbarkeit) benennen,
  \item zu einer klar formulierten Entscheidung f\"uhren \textendash{} nicht blo{\ss} zu einer Beschreibung.
\end{itemize}

\begin{infobox}[Nutzung im Coaching]
Vergleichen Sie Ihre eigene Antwort \emph{vor} der Musterl\"osung. Wo weicht Ihre Argumentation ab? Ist die Abweichung eine andere Annahme \textendash{} oder ein Denkfehler? Genau dort entsteht der Lerneffekt.
\end{infobox}

\clearpage

% =========================================================================
% L1 LOESUNGEN
% =========================================================================
\section*{Niveau L1 \textendash{} Wissen \& Reproduktion}
\addcontentsline{toc}{section}{L1 -- Wissen}

% --- AUFGABE 1 -----------------------------------------------------------
% LERNPLATT_HILFSVIDEOS
\section*{Hilfsvideos zu diesem Tag}
\addcontentsline{toc}{section}{Hilfsvideos zu diesem Tag}
\thispagestyle{plain}

\noindent Die folgenden YouTube-Erkl\"arvideos vermitteln die Inhalte, die Sie zur Bearbeitung der Aufgaben ben\"otigen. Klicken Sie auf den Titel, um das Video direkt zu \"offnen; die vollst\"andige URL steht in Klammern.

\begin{description}[style=nextline,leftmargin=18pt,labelindent=0pt,itemsep=8pt]
  \item[1.~Microservices — warum, wann und mit welchen Grenzen] \href{https://youtu.be/ZXfLy229GLs}{\textcolor{lpAccent}{https://youtu.be/ZXfLy229GLs}}\\[2pt]
  {\footnotesize\color{lpGray}(URL: \nolinkurl{https://youtu.be/ZXfLy229GLs})}
  \item[2.~Kubernetes — Container-Orchestrierung in der Praxis] \href{https://youtu.be/SFVg7fDdvLE}{\textcolor{lpAccent}{https://youtu.be/SFVg7fDdvLE}}\\[2pt]
  {\footnotesize\color{lpGray}(URL: \nolinkurl{https://youtu.be/SFVg7fDdvLE})}
  \item[3.~Test-Driven Development — Qualität als Entscheidungsinstrument] \href{https://youtu.be/71nLhdZuMk0}{\textcolor{lpAccent}{https://youtu.be/71nLhdZuMk0}}\\[2pt]
  {\footnotesize\color{lpGray}(URL: \nolinkurl{https://youtu.be/71nLhdZuMk0})}
  \item[4.~DevOps-Operating-Model — „you build it, you run it” und Plattform-Teams] \href{https://youtu.be/ak3oE0I6cXU}{\textcolor{lpAccent}{https://youtu.be/ak3oE0I6cXU}}\\[2pt]
  {\footnotesize\color{lpGray}(URL: \nolinkurl{https://youtu.be/ak3oE0I6cXU})}
\end{description}
\vspace{0.6em}
\noindent{\footnotesize\color{lpGray}\textit{Tipp:} \"Offnen Sie das Video, lassen Sie es im Hintergrund laufen und arbeiten Sie parallel an der Aufgabe.}

\clearpage

\aufgabenkopf{1}{Microservices \textendash{} Eignung, Risiken und Servicegrenze}{\nivLi}{8\,min}

\ytquery{Microservices Vorteile Risiken Servicegrenze Business Capability erklaert}

\begin{loesungbox}
\textbf{1. Drei Nutzenversprechen:}
\begin{itemize}
  \item \emph{Unabh\"angige Deployments:} Teams releasen ihren Service ohne Koordination mit allen anderen \textendash{} k\"urzere Lead Time, niedrigere Stillstandskosten.
  \item \emph{Gezielte Skalierung:} Nur der Service, der unter Last steht, wird skaliert (z.\,B.\ \emph{Checkout} zur Peak-Saison). Kein Skalieren ``auf Verdacht''.
  \item \emph{Klare Verantwortungsgrenzen:} Ein Service, ein Team, ein Owner. ``Verantwortung verschwindet nicht im Korridor.''
\end{itemize}

\smallskip
\textbf{2. Drei Risiken:}
\begin{itemize}
  \item \emph{Komplexit\"at:} Viele bewegliche Teile, Netzwerk als unzuverl\"assige Komponente, Latenzen, Versionierung der Schnittstellen.
  \item \emph{Verteilte Fehler:} Ein Fehler wandert \"uber Service-Grenzen; Debugging ist ohne saubere Korrelations-IDs sehr teuer.
  \item \emph{Observability-Bedarf:} Logs, Metriken, Traces zentral und konsistent \textendash{} sonst dauert die Root-Cause-Analyse bei jedem Incident Stunden.
\end{itemize}

\smallskip
\textbf{3. Gute Servicegrenze:} Eine Servicegrenze ist \emph{gut}, wenn sie entlang einer fachlichen Capability verl\"auft (z.\,B.\ \emph{Inventory}, \emph{Billing}) und \emph{nicht} entlang von Datenbanktabellen. Tabellen sind Implementierungsdetail; eine Capability ist die kleinste Einheit gesch\"aftlicher Verantwortung \textendash{} sie \"uberlebt Refactorings und definiert sauberen Owner-Schnitt.
\end{loesungbox}

\clearpage

% --- AUFGABE 2 -----------------------------------------------------------
\aufgabenkopf{2}{Container \& Kubernetes \textendash{} Kernbegriffe}{\nivLi}{8\,min}

\ytquery{Kubernetes Grundlagen Container Orchestrierung Blue Green Canary einfach erklaert}

\begin{loesungbox}
\begin{tabularx}{\linewidth}{@{}p{3.6cm}X@{}}
\toprule
\textbf{Begriff} & \textbf{Erkl\"arung} \\
\midrule
Container & Verpackte Anwendung inklusive aller Abh\"angigkeiten \textendash{} l\"auft auf jedem kompatiblen Host gleich. ``Build once, run anywhere.'' \\
Orchestrierung & Plattform-Layer, der Container automatisch startet, neu plant, vernetzt, skaliert und ausf\"allige neu hochzieht (Kubernetes als typische Implementierung). \\
Self-Healing & Wenn ein Container abst\"urzt, startet die Orchestrierung ihn automatisch neu \textendash{} ohne menschliches Eingreifen. \\
Rollback & R\"uckkehr zu einer vorherigen, bekannten guten Version, wenn ein neues Release Fehler zeigt. \\
Service-Discovery & Services finden sich \"uber logische Namen, nicht \"uber IP-Adressen \textendash{} \"Anderungen am Cluster brechen Aufrufer nicht. \\
Blue/Green-Deployment & Zwei vollst\"andige Versionen laufen parallel; der Verkehr wird in einem Schritt von ``blau'' auf ``gr\"un'' umgeschaltet. \\
Canary-Release & Eine kleine Menge des Verkehrs (z.\,B.\ 5\,\%) geht zur neuen Version; bei stabilen Kennzahlen wird der Anteil schrittweise erh\"oht. \\
Ressourcen\-limits & Pro Container vereinbarte Mindest- (Requests) und H\"ochstwerte (Limits) f\"ur CPU/RAM \textendash{} verhindert, dass ein Service den Cluster ``frisst''. \\
\bottomrule
\end{tabularx}

\smallskip
\textbf{Zwei kritischste:} \emph{Rollback} und \emph{Ressourcenlimits}. Ohne ge\"ubten Rollback wird jeder Incident zur Improvisation; ohne Limits bringt ein einzelner fehlerhafter Service die Plattform in den Ressourcen-Tod. Beide kosten in der Praxis am schnellsten am meisten Geld.
\end{loesungbox}

\clearpage

% --- AUFGABE 3 -----------------------------------------------------------
\aufgabenkopf{3}{Testpyramide \& Quality Gates}{\nivLi}{10\,min}

\ytquery{Testpyramide TDD Red Green Refactor Quality Gates erklaert}

\begin{loesungbox}
\begin{tabularx}{\linewidth}{@{}p{3.4cm}|X|X|X@{}}
\toprule
\textbf{Eigenschaft} & \textbf{Unit-Test} & \textbf{Integrationstest} & \textbf{E2E-Test} \\
\midrule
Was wird getestet? & Einzelne Funktion / Klasse isoliert. & Zusammenspiel von 2\,--\,3 Komponenten (z.\,B.\ Service + DB). & Ganzer Gesch\"aftsfluss \"uber mehrere Services. \\
\midrule
Geschwindigkeit & Sehr schnell (Millisekunden). & Mittel (Sekunden). & Langsam (Minuten). \\
\midrule
Kosten / Fragilit\"at & G\"unstig, stabil. & Mittel, m\"a{\ss}ig fragil. & Teuer, sehr fragil (UI, Netzwerk, Timing). \\
\midrule
Anteil in Pyramide & \textbf{Viele} ($\sim$\,70\,\%). & Mittel ($\sim$\,20\,\%). & \textbf{Wenige} ($\sim$\,10\,\%). \\
\midrule
Typischer Fehlerfund & Logikfehler in einer Methode. & Schnittstelle / Datenformat falsch. & Geh\"art zusammen, aber im Gesamtfluss bricht etwas. \\
\bottomrule
\end{tabularx}

\smallskip
\textbf{Quality Gates} sind verbindliche Pr\"ufungen in der Pipeline, die jeder Build durchlaufen muss, bevor er weiter darf. Drei nicht verhandelbare Gates vor einem Produktiv-Release:

\begin{enumerate}
  \item \emph{Alle Unit- und Integrationstests gr\"un.}
  \item \emph{Static Code- und Security-Scan ohne kritische Findings} (z.\,B.\ keine Secrets im Code, keine bekannten CVEs in Dependencies).
  \item \emph{Smoke Test nach Deployment} (mindestens ein ``Happy Path''-Aufruf gegen die neue Version).
\end{enumerate}
\end{loesungbox}

\clearpage

% =========================================================================
% L2 LOESUNGEN
% =========================================================================
\section*{Niveau L2 \textendash{} Anwendung \& Transfer}
\addcontentsline{toc}{section}{L2 -- Anwendung}

% --- AUFGABE 4 -----------------------------------------------------------
\aufgabenkopf{4}{Service-Schnitt aus Prozessbeschreibung \textendash{} Order-to-Cash}{\nivLii}{25\,min}

\ytquery{Microservices Service Schnitt Order to Cash Domain Driven Design Capability}

\begin{loesungbox}
\textbf{1. Sechs Service-Kandidaten:}
\begin{itemize}
  \item \emph{Order Service} \textendash{} Bestellungen annehmen, validieren, im Lifecycle f\"uhren.
  \item \emph{Inventory Service} \textendash{} Verf\"ugbarkeit pr\"ufen, Reservierungen, Best\"ande.
  \item \emph{Pricing/Quote Service} \textendash{} Preise und Angebote berechnen.
  \item \emph{Shipping Service} \textendash{} Versandauftr\"age, Tracking.
  \item \emph{Billing Service} \textendash{} Rechnungen erzeugen, Versand, Mahnwesen.
  \item \emph{Payments Service} \textendash{} Zahlungen abwickeln, R\"uckbuchungen.
  \item \emph{Complaints Service} \textendash{} Reklamationen, Gutschriften, Eskalation.
\end{itemize}

\smallskip
\textbf{2. Drei Service-Steckbriefe:}

\smallskip
\noindent\textbf{(a) Order Service.} \emph{Zweck:} Lebenszyklus einer Bestellung halten (angelegt \textrightarrow{} best\"atigt \textrightarrow{} verschickt \textrightarrow{} abgeschlossen). \emph{Input:} Kunde, Artikelpositionen, Lieferadresse, gew\"unschter Liefertermin. \emph{Output:} Bestell-ID, Status, best\"atigter Liefertermin, Gesamtsumme. \emph{Owner:} Produktteam Order. \emph{Kritischster Fehler:} Bestellung in inkonsistentem Zustand (best\"atigt, aber Inventory nicht reserviert) \textendash{} Doppelt-Versand oder Lieferversprechen ohne Deckung.

\smallskip
\noindent\textbf{(b) Payments Service.} \emph{Zweck:} Zahlungen f\"ur best\"atigte Bestellungen verarbeiten und buchen. \emph{Input:} Bestell-ID, Betrag, Zahlungsmittel, IBAN/Karteninfo (via Token). \emph{Output:} Payment-Status, Buchungs-ID, Beleg-Referenz, Fehlerursache. \emph{Owner:} Produktteam Payments + Compliance. \emph{Kritischster Fehler:} Doppelbuchung oder fehlgeschlagene Buchung ohne Storno \textendash{} direkte finanzielle und reputative Folgen.

\smallskip
\noindent\textbf{(c) Inventory Service.} \emph{Zweck:} Best\"ande, Reservierungen und Verf\"ugbarkeitspr\"ufungen. \emph{Input:} Artikel-ID, Menge, Standort, Reservierungszeitraum. \emph{Output:} Verf\"ugbarkeit (ja/nein/eingeschr\"ankt), Reservierungs-ID, voraussichtliches Wiederverf\"ugbarkeits-Datum, Sicherheitsbestand-Status. \emph{Owner:} Produktteam Inventory. \emph{Kritischster Fehler:} Verkauf von tats\"achlich nicht vorhandenem Bestand (``Overselling'').

\smallskip
\textbf{3. Riskanteste Grenze + Guardrail:} \emph{Order \textendash{} Inventory.} Hier kollidieren zwei Verantwortlichkeiten: Order \emph{verspricht}, Inventory \emph{deckt}. Datenhoheit ist h\"aufig unklar (Reservierung als Order- oder Inventory-Konzept?). \emph{Guardrail:} \textbf{Anti-Corruption-Layer + idempotente Reservierungs-API}: Order ruft Inventory mit eindeutiger Operations-ID; Inventory antwortet entweder ``reserviert'' (mit TTL) oder ``nicht verf\"ugbar''. Doppelaufrufe wirken sich \emph{nicht} mehrfach aus. Ein \emph{gemeinsames Datenmodell} ist verboten \textendash{} jede Seite h\"alt ihre eigene Sicht.

\smallskip
\emph{Andere Schnitte sind legitim} (z.\,B.\ Trennung Pricing in \emph{Quote} + \emph{Pricing}, oder \emph{Notifications} als eigener Service). Entscheidend: Capability-Logik, klarer Owner, niedrige Kopplung.
\end{loesungbox}

\clearpage

% --- AUFGABE 5 -----------------------------------------------------------
\aufgabenkopf{5}{Release-Flow \& Testpyramide als Sicherheitsnetz}{\nivLii}{25\,min}

\ytquery{Release Flow Quality Gates Canary Feature Toggle SLO Beispiel}

\begin{loesungbox}
\textbf{1. Minimaler Release-Flow (7 Schritte):}
\begin{enumerate}
  \item \emph{Commit + PR} \textendash{} Code-Review durch zweite Person.
  \item \emph{Build + Unit/Integrationstests} \textendash{} \textbf{Quality Gate 1}: Tests gr\"un, Coverage $\geq$\,80\,\%.
  \item \emph{Static \& Security Scan} \textendash{} Dependency-Check, Secret-Scan, License-Check.
  \item \emph{Stage-Deployment + Smoke Test} \textendash{} Happy Path und einen Fehlerpfad pr\"ufen.
  \item \emph{Approval} \textendash{} \textbf{Quality Gate 2}: Service-Owner gibt explizit frei (bei kritischen Services zus\"atzlich Security).
  \item \emph{Canary auf Prod} \textendash{} 5\,\% Verkehr, 30 Minuten Beobachtung (Fehlerquote, Latenz, Logs).
  \item \emph{Full Rollout} bei stabilen Kennzahlen, sonst automatischer Rollback.
\end{enumerate}

\smallskip
\textbf{2. Testpyramide f\"ur drei Services:}

\smallskip
\begin{tabularx}{\linewidth}{@{}lXXX@{}}
\toprule
\textbf{Service} & \textbf{Unit} & \textbf{Integration} & \textbf{E2E} \\
\midrule
Order   & $\sim$\,120 (Geschaeftsregeln, Statuslogik) & $\sim$\,25 (DB, Inventory-Aufruf)  & 2 (Happy Path, Fraud-Pfad) \\
Billing & $\sim$\,90 (Berechnung, Steuerlogik)        & $\sim$\,20 (PDF-Erzeugung, Versand) & 1 (Rechnung an Kunde) \\
Payments & $\sim$\,140 (PSP-Adapter, Storno-Logik)    & $\sim$\,30 (Zahlungsprovider, Buchung) & 2 (Erfolgsfall, Rueckbuchung) \\
\bottomrule
\end{tabularx}

\smallskip
\textbf{3. Zwei Betriebskennzahlen (SLO-orientiert):}
\begin{itemize}
  \item \emph{Verf\"ugbarkeit Order-API:} 99,9\,\% pro Monat. \emph{Owner:} Service-Owner Order. \emph{Mess-fenster:} rollierende 30 Tage.
  \item \emph{Fehlerquote Payments:} $<$\,0,5\,\% der Buchungsversuche. \emph{Owner:} Service-Owner Payments + Plattformteam (Alerts). \emph{Mess-fenster:} rollierende 7 Tage.
\end{itemize}

\smallskip
\textbf{4. MVP-Risikoplan (Rest-Risiko + Kompensation):}
\begin{itemize}
  \item \emph{Bewusster Verzicht:} weniger E2E-Tests (nur 1\,--\,2 Kernpfade pro Service).
  \item \emph{Kompensation 1:} \emph{Canary-Release} mit harten Alerting-Schwellen (Fehlerquote, p95-Latenz) \textendash{} kein full rollout, solange Kennzahlen unklar.
  \item \emph{Kompensation 2:} \emph{Feature Toggles} f\"ur riskante \"Anderungen \textendash{} \"Anderungen lassen sich ohne Deployment ausschalten.
  \item \emph{Kompensation 3:} \emph{Manuelle Stichprobe} f\"ur Zahlungen mit Betr\"agen $\geq$\,1.000\,\euro{} in den ersten zwei Wochen nach gr\"o{\ss}eren Releases.
\end{itemize}
\end{loesungbox}

\clearpage

% --- AUFGABE 6 -----------------------------------------------------------
\aufgabenkopf{6}{Fallstudie ``InsureFast'' \textendash{} Servicegrenzen \& Betriebskonzept}{\nivLii}{30\,min}

\ytquery{Microservices Versicherung Schadenmeldung Kubernetes Skalierung Beispiel}

\begin{loesungbox}
\textbf{1. Servicegrenzen (sieben Services):}

\smallskip
\begin{itemize}
  \item \emph{Claim Intake} \textendash{} Eingang aller Schadenmeldungen (App, Telefon, Mail). Validierung, Vorqualifizierung.
  \item \emph{Policy/Customer} \textendash{} Vertrags- und Kundendaten (Datenhoheit Kunde).
  \item \emph{Fraud Check} \textendash{} Betrugsindikatoren, automatisierte Pr\"ufung, Flagging.
  \item \emph{Claim Assessment} \textendash{} Sachpr\"ufung, Entscheidung, Begr\"undung (auditrelevant!).
  \item \emph{Payout} \textendash{} Auszahlung, Verkn\"upfung mit Finance.
  \item \emph{Notifications} \textendash{} Mail, Push, SMS an Kunden + interne Stakeholder.
  \item \emph{Document Management} \textendash{} Beweisdokumente, Fotos, Gutachten, Archivierung.
\end{itemize}

\smallskip
\textbf{Begr\"undung einer nicht-trivialen Grenze:} \emph{Fraud Check} als eigener Service (nicht in \emph{Claim Assessment}). Grund: Fraud Check ist regelm\"a{\ss}ig algorithmisch-dynamisch (Modelle, Schwellenwerte), wird h\"aufiger ge\"andert und braucht eigene Datenhoheit (Score-Historie). Eingebettet in Assessment w\"urde jede Modell\"anderung den gesamten Pr\"ufprozess deployen lassen.

\smallskip
\textbf{2. Kubernetes-Betriebskonzept:}
\begin{itemize}
  \item \emph{Rollout/Rollback:} \textbf{Canary} f\"ur Claim Intake und Assessment (Verkehrsschritte 5\,\%/25\,\%/100\,\%); Rollback automatisch ausgel\"ost, wenn Fehlerquote $>$\,1\,\% \"uber 5 Minuten. Pflicht: Rollback innerhalb von 10 Minuten ge\"ubt.
  \item \emph{Skalierungsregel:} ``Wenn die Queue-L\"ange in \emph{Claim Intake} $>$\,200 oder die CPU-Auslastung $>$\,70\,\% f\"ur 3 Minuten, dann +2 Replicas, bis max.\ 12.''
  \item \emph{Konfiguration / Secrets:} Konfig in ConfigMaps; Secrets im Vault (nie im Image); Zugriff role-based; Rotation alle 90 Tage.
  \item \emph{Monitoring / Alerting:} Zentraler Log-/Metric-/Trace-Standard mit Korrelations-ID pro Schadenfall. Alerts auf Fehlerquote, Latenz (p95), Queue-L\"ange, fehlende Heartbeats; Eskalation an Service-Owner.
  \item \emph{Ressourcenlimits:} Pro Container harte Limits (z.\,B.\ 500\,m CPU, 512\,MiB RAM) und Requests bei 70\,\% der Limits. Keine ``unbegrenzten'' Pods.
\end{itemize}
\end{loesungbox}

\clearpage

% --- AUFGABE 7 -----------------------------------------------------------
\aufgabenkopf{7}{Fallstudie ``InsureFast'' \textendash{} Teststrategie \& Operating Model}{\nivLii}{30\,min}

\ytquery{Microservices Operating Model RACI Plattform Team Incident Management}

\begin{loesungbox}
\textbf{1. Teststrategie (TDD):}
\begin{itemize}
  \item \emph{Unit-Tests} f\"ur Kernlogik: \emph{Claim Assessment} (Entscheidungsregeln, Begr\"undungslogik), \emph{Fraud Check} (Score-Berechnung), \emph{Payout} (Storno-/Korrekturlogik). TDD-Zyklus \emph{Red--Green--Refactor} verpflichtend f\"ur regulatorische Logik.
  \item \emph{Integrationstests} f\"ur Service-APIs: Claim Intake \textrightarrow{} Assessment, Assessment \textrightarrow{} Payout, alle Services \textrightarrow{} Document Management.
  \item \emph{E2E-Tests} nur f\"ur \textbf{zwei Kernpfade}: \emph{Happy Path} (Schaden gemeldet \textrightarrow{} gepr\"uft \textrightarrow{} ausgezahlt) und \emph{Fraud-Flag-Path} (Schaden gemeldet \textrightarrow{} Fraud Check schl\"agt an \textrightarrow{} manueller Review).
  \item \emph{Quality Gates:} (1) Alle Tests gr\"un + Security-Scan ohne kritische Findings; (2) Smoke Test nach Deploy, Rollback-Fallback geplant.
\end{itemize}

\smallskip
\textbf{2. Operating Model:}

\smallskip
\noindent\emph{Teamzuschnitt:} 5\,--\,6 Produktteams, jedes Team besitzt 1\,--\,2 Services (``you build it, you run it''). 1 Plattformteam (CI/CD-Templates, Observability, Kubernetes-Basis). 1 Security/Compliance-Team (Baselines, Ausnahmen).

\smallskip
\noindent\emph{RACI Schnittstellen\"anderung:}

\smallskip
\begin{tabularx}{\linewidth}{@{}lcccc@{}}
\toprule
 & \textbf{Service-Owner (anrufend)} & \textbf{Service-Owner (gerufen)} & \textbf{Plattform} & \textbf{Security} \\
\midrule
\"Anderung initiieren     & R / A & C & I & I \\
Vertragspr\"ufung         & C     & R & C & I \\
Auswirkung absch\"atzen   & C     & A & R & C \\
Freigabe                 & I     & A & C & R / C \\
Rollback-Plan            & R     & R & A & I \\
\bottomrule
\end{tabularx}

\smallskip
\noindent\emph{Incident-Handling:} On-call beim Service-Owner; bei Plattformproblemen Eskalation an Plattformteam; Security wird zugeschaltet bei Datenleck-Verdacht. Postmortem innerhalb von 5 Arbeitstagen, blameless.

\smallskip
\noindent\emph{Release-Freigabepfad:} Standard-Release \textrightarrow{} Service-Owner gen\"ugt. Kritische \"Anderung (Schnittstelle, Datenmodell, Security) \textrightarrow{} Service-Owner + Plattform + Security.

\smallskip
\textbf{3. Zwei Entscheidungen unter Unsicherheit:}
\begin{itemize}
  \item \emph{Datenstrategie:} \textbf{DB-bleibt + Anti-Corruption-Layer} (kein Big Bang). Begr\"undung: 12 Wochen reichen nicht f\"ur Migration unter Audit-Druck; Risiko: Migrations-Bug w\"are existenzkritisch. \emph{Fr\"uhwarnsignal:} Wenn die Integrationsfehler \"uber den Anti-Corruption-Layer $>$\,5\,\% steigen oder die Lead-Time pro \"Anderung sich verdoppelt \textrightarrow{} Migration einplanen.
  \item \emph{Release-Frequenz:} \textbf{w\"ochentlich f\"ur Nicht-Kernservices, zweiw\"ochentlich f\"ur Kern} (Assessment, Payout). \emph{Fr\"uhwarnsignal:} Change Failure Rate $>$\,15\,\% im Kern \textrightarrow{} Frequenz reduzieren oder Quality Gates verst\"arken.
\end{itemize}
\end{loesungbox}

\clearpage

% --- AUFGABE 8 -----------------------------------------------------------
\aufgabenkopf{8}{Anti-Muster ``Microservices-Spaghetti'' erkennen und korrigieren}{\nivLii}{25\,min}

\ytquery{Microservices Antipattern verteilter Monolith Observability erklaert}

\begin{loesungbox}
\textbf{A \textendash{} Service-Wildwuchs.}
\begin{itemize}
  \item \emph{Verletztes Prinzip:} ``Servicegrenze entlang Capabilities'' und ``Wirtschaftlichkeit''. 47 Services f\"ur 60 Entwickler bedeuten weniger als 1,3 Entwickler pro Service \textendash{} unten \"uber Owner und Run-Disziplin.
  \item \emph{Korrektur:} Konsolidieren entlang Capabilities. Faustregel: ein Service braucht mindestens \emph{drei Personen}, die ihn ge- und betreiben k\"onnen, sonst zu Capability verschmelzen. Kundendaten in einer Hoheit \textendash{} \emph{Customer Service} \textendash{} statt vierfach.
\end{itemize}

\smallskip
\textbf{B \textendash{} Verteilter Monolith.}
\begin{itemize}
  \item \emph{Verletztes Prinzip:} ``Lose Kopplung'' und ``unabh\"angige Deployments''. Wenn jede \"Anderung mehrere Releases erzwingt, sind die Services nur formal getrennt.
  \item \emph{Korrektur:} Synchrone Aufrufe zu Events / asynchroner Kommunikation umbauen, wo m\"oglich. API-Versionierung erzwingen (alte Version mindestens eine Iteration weiter laufen lassen). Schnittstellen explizit dokumentieren (OpenAPI/AsyncAPI) und Vertragsbruch-Tests f\"ahren.
\end{itemize}

\smallskip
\textbf{C \textendash{} Observability als Nachgedanke.}
\begin{itemize}
  \item \emph{Verletztes Prinzip:} Betriebsdisziplin, ``you build it, you run it''.
  \item \emph{Korrektur:} Zentrale Log-/Metric-/Trace-Standards als Pflicht (Plattformteam stellt sie bereit, nicht jedes Team baut sie nach). \emph{Korrelations-ID} pro Gesch\"aftsfall durch alle Services. Mindest-Dashboard pro Service: Verkehr, Fehlerquote, Latenz, Saturation (``USE/RED-Pattern''). Ziel-MTTR explizit.
\end{itemize}

\smallskip
\textbf{D \textendash{} Rollback-Atrophie.}
\begin{itemize}
  \item \emph{Verletztes Prinzip:} Betriebsdisziplin und Reliability-Engineering. Ein nie ge\"ubter Rollback ist im Incident-Fall \emph{kein} Rollback.
  \item \emph{Korrektur:} Monatliche Rollback-\"Ubung pro Team (Chaos-Light), Rollback-Pfad in jeder Pipeline implementiert und getestet, Rollback-Dauer als SLO f\"ur das Plattformteam (z.\,B.\ $\leq$\,10\,Minuten). Postmortems erw\"ahnen explizit, ob Rollback m\"oglich gewesen w\"are.
\end{itemize}
\end{loesungbox}

\clearpage

% =========================================================================
% L3 LOESUNGEN
% =========================================================================
\section*{Niveau L3 \textendash{} Management \& Entscheidungsarchitektur}
\addcontentsline{toc}{section}{L3 -- Management}

% --- AUFGABE 9 -----------------------------------------------------------
\aufgabenkopf{9}{Reliability- \& Risk-Framework f\"ur eine wachsende Plattform}{\nivLiii}{35\,min}

\ytquery{SLO Error Budget Change Failure Rate MTTR SRE Framework Senior}

\begin{loesungbox}
\textbf{1. SLO-Set (drei SLOs f\"ur kritische Services):}

\smallskip
\begin{tabularx}{\linewidth}{@{}p{3.6cm}lX@{}}
\toprule
\textbf{SLO} & \textbf{Schwelle} & \textbf{Begr\"undung} \\
\midrule
Verf\"ugbarkeit & 99,9\,\% / Monat & Branchen\"ublich f\"ur Kundenkernprozesse; entspricht ca.\ 43 Minuten Ausfall pro Monat. \\
Latenz p95     & $\leq$\,300\,ms  & Schwelle, ab der Nutzerwahrnehmung sp\"urbar leidet (User Research). \\
Fehlerquote    & $\leq$\,0,5\,\%  & Niedrige Schwelle, weil Fehler in der Kernlogik direkt in Reklamationen und Audit-Findings durchschlagen. \\
\bottomrule
\end{tabularx}

\smallskip
\textbf{2. Error Budgets:} Wenn die Verf\"ugbarkeit 99,9\,\% / Monat als SLO definiert ist, betr\"agt das \emph{Error Budget} 0,1\,\% / Monat ($\sim$\,43 Minuten). Solange das Budget \emph{nicht} aufgebraucht ist, d\"urfen Teams releasen. \emph{Wird das Budget aufgebraucht, gilt:}
\begin{itemize}
  \item Release-Stop f\"ur den betroffenen Service \textendash{} ausgenommen sicherheitskritische Hotfixes.
  \item Plattform-Review innerhalb von 5 Arbeitstagen mit dem Service-Owner: Ursache, Ma{\ss}nahme, Lernen.
  \item Erst nach 14 Tagen ohne weitere Verletzung wird der Release-Stop aufgehoben.
\end{itemize}
\emph{Sinn:} Speed und Reliability werden \emph{verhandelbar} \textendash{} kein blo{\ss}er Appell.

\smallskip
\textbf{3. Risk-Framework (f\"unf Risiken):}

\smallskip
\begin{tabularx}{\linewidth}{@{}p{4.0cm}ccX@{}}
\toprule
\textbf{Risiko} & \textbf{Impact} & \textbf{Whk.} & \textbf{Pr\"avention} \\
\midrule
Servicegrenze falsch geschnitten & 5 & 3 & Capability-Mapping vor Schnitt, Architekturreview f\"ur irreversible Entscheidungen. \\
Secret im Code & 5 & 3 & Secret-Scanning als Quality Gate; Vault verpflichtend. \\
Rollback nicht ge\"ubt & 4 & 4 & Monatliche Rollback-\"Ubung, Rollback-Dauer als SLO. \\
Observability-L\"ucken & 4 & 4 & Mindest-Dashboard verpflichtend, Korrelations-ID Standard. \\
Tooling-Lock-in & 3 & 3 & Offene Standards (OpenAPI, BPMN-XML, Helm), Export-Anforderung in Tool-Auswahl. \\
\bottomrule
\end{tabularx}

\smallskip
\textbf{4. Anti-Gaming-Mechanik:}
\begin{itemize}
  \item \emph{Kopplung 1:} Change Failure Rate gilt nur dann als ``erf\"ullt'', wenn gleichzeitig der \emph{Service-Health-Index} (interner Composite-Wert aus Fehlern + Eskalationen + Retouren) sich nicht verschlechtert. So bringt das Klein-Halten von ``Changes'' keinen scheinbaren Erfolg.
  \item \emph{Kopplung 2:} Postmortem-Pflicht f\"ur jeden \emph{Kundenwahrnehmbaren} Incident; ein nicht durchgef\"uhrter Postmortem z\"ahlt als Change Failure mit doppelter Gewichtung. So lohnt es sich nicht, Incidents zu vertuschen.
\end{itemize}
\end{loesungbox}

\clearpage

% --- AUFGABE 10 ----------------------------------------------------------
\aufgabenkopf{10}{Governance leicht, aber wirksam}{\nivLiii}{30\,min}

\ytquery{Architecture Governance Principles Senior Lightweight Service Standard}

\begin{loesungbox}
\textbf{1. Acht Architekturprinzipien:}
\begin{enumerate}
  \item Servicegrenzen folgen \emph{Business Capabilities}, nicht Datenbanktabellen oder Org-Charts.
  \item Jeder Service hat \emph{einen} Owner; Ownership ist an die Funktion gekoppelt, nicht an Personen.
  \item Kein Secret im Code; Credentials kommen ausschlie{\ss}lich aus dem Vault.
  \item Jede \"offentliche Schnittstelle ist als OpenAPI/AsyncAPI dokumentiert und versioniert; Vertragsbr\"uche werden in der Pipeline erkannt.
  \item Schreibende Operationen sind \emph{idempotent} (Operations-ID), Konsumenten d\"urfen sicher wiederholen.
  \item Jeder Service stellt ein Mindest-Dashboard (Verkehr, Fehlerquote, Latenz, Saturation) bereit und schreibt strukturierte Logs mit Korrelations-ID.
  \item Releases gehen \"uber Canary oder Blue/Green; Rollback ist innerhalb von 10 Minuten m\"oglich und monatlich ge\"ubt.
  \item Datenhoheit ist eindeutig: ein Service ``besitzt'' eine Capability, Daten anderer Services werden \"uber Schnittstellen oder Events bezogen \textendash{} nicht via geteilter Datenbank.
\end{enumerate}

\smallskip
\textbf{2. Schnittstellenstandards (zwei):}
\begin{itemize}
  \item \emph{OpenAPI / AsyncAPI als Single Source of Truth} f\"ur synchrone bzw.\ asynchrone Schnittstellen \textendash{} jeder Service ver\"offentlicht seine Schemas im Service-Katalog.
  \item \emph{Event-Schema-Registry} (z.\,B.\ Avro/JSON-Schema) mit Versionierung und R\"uckw\"artskompatibilit\"ats-Pr\"ufung als Quality Gate.
\end{itemize}

\smallskip
\textbf{3. Review-Regel:}
\begin{itemize}
  \item \emph{Architektur-Review verpflichtend bei:} neuer Service, neue Schnittstelle nach au{\ss}en, Datenhoheits-\"Anderung, Plattformwechsel, Sicherheits-Baseline-Ausnahmen.
  \item \emph{Nicht erforderlich bei:} interne Refactorings, neue interne Endpoints, Test-/Toolwechsel innerhalb des Teams.
\end{itemize}
\emph{Faustregel:} Review nur bei \emph{teuren oder irreversiblen} Entscheidungen \textendash{} alles andere kostet mehr Zeit als es schafft.

\smallskip
\textbf{4. Eskalationspfad (vier Stufen):}
\begin{enumerate}
  \item Service-Owner und Plattform-Lead versuchen Einigung (max.\ 2 Arbeitstage).
  \item Bei Patt: Architekturboard (Plattform-Lead + zwei rotierende Senior-Engineers + Security).
  \item Bei strategischem Dissens: Head of Platform / Head of Product.
  \item Bei strategischer Verantwortung: CTO. Entscheidungen sind dokumentiert (ADR \textendash{} Architecture Decision Record).
\end{enumerate}
\end{loesungbox}

\clearpage

% --- AUFGABE 11 ----------------------------------------------------------
\aufgabenkopf{11}{Transferprojekt \textendash{} Digitale Prozessarchitektur als Betriebs- und Entscheidungsarchitektur}{\nivLiii}{45\,min}

\ytquery{Microservices Operating Model Decision Architecture Senior Platform CIO}

\begin{loesungbox}
\emph{Musterl\"osung als Entscheidungsarchitektur \textendash{} andere Konfigurationen sind m\"oglich, solange sie als Architektur (nicht als Liste) gedacht sind.}

\smallskip
\textbf{1. Top-10-Entscheidungsarchitektur:}

\smallskip
\begin{tabularx}{\linewidth}{@{}p{3.8cm}p{3.6cm}X@{}}
\toprule
\textbf{Entscheidung} & \textbf{Wer entscheidet?} & \textbf{Optionen / Kriterien} \\
\midrule
Servicegrenzen          & Architekturboard + Service-Owner & Capability-Schnitt; max.\ 20 Services in Phase 1. \\
Datenhoheit             & Architekturboard                & Eine Capability = eine Datenhoheit; kein Shared-DB. \\
API-/Event-Standards    & Plattform-Lead                  & OpenAPI/AsyncAPI verbindlich; Schema-Registry. \\
SLOs                    & Service-Owner + Plattform       & Standard-Set (Verf./Lat./Err) je Klasse. \\
Release-Policies        & Plattform-Lead + Security       & Standard-Releases via Canary; Kern-Release zweiw\"ochentlich. \\
Security-Baselines      & CISO + Plattform                & Vault, Least Privilege, SBOM, Secret-Scan. \\
Observability-Standard  & Plattform                        & Zentraler Log-/Metric-/Trace-Stack; Korrelations-ID Pflicht. \\
Incident-Governance     & Head of Platform                 & On-call beim Service-Owner; Eskalation an Plattform. \\
Ausnahmeprozess         & Architekturboard                 & Schriftliche Ausnahme, Ablaufdatum, Reviewer. \\
Toolchain-Standard      & CTO                              & 1\,--\,2 Tools pro Kategorie; Export-Pflicht. \\
\bottomrule
\end{tabularx}

\smallskip
\textbf{2. Operating Model \& RACI:}
\begin{itemize}
  \item \emph{Produktteams (Build + Run):} Service-Owner, Tag-Betrieb, Releases, On-call, Postmortems.
  \item \emph{Plattformteam:} CI/CD-Templates, Observability, Kubernetes-Basis, Rollback-Mechanik, Schema-Registry.
  \item \emph{Security/Compliance:} Baselines, Audit-Beweise, Ausnahmen, SBOM-Pflicht.
  \item \emph{Architekturboard:} nur teure/irreversible Entscheidungen; max.\ 1\,h pro Sitzung, alle 2 Wochen.
  \item \emph{Eskalation bei Speed vs.\ Risk:} Service-Owner \textrightarrow{} Plattform-Lead \textrightarrow{} Head of Platform \textrightarrow{} CTO.
\end{itemize}

\smallskip
\textbf{3. Reliability- \& Risk-Framework:}
\begin{itemize}
  \item SLOs pro Service (Standard-Set), Error Budgets, Change Failure Rate ($\leq$\,10\,\%), MTTR ($\leq$\,60\,min).
  \item Anti-Gaming-Kopplung: Kosten-/Geschwindigkeits-KPIs gelten nur, wenn \emph{Service-Health-Index} nicht degradiert (s.\ Aufgabe 9).
\end{itemize}

\smallskip
\textbf{4. Governance leicht, aber wirksam:} max.\ 8 Architekturprinzipien (s.\ Aufgabe 10). Standard-Templates f\"ur CI/CD, Helm/Deployment, Mindest-Dashboard, Postmortem.

\smallskip
\textbf{5. Roadmap (16 Wochen):}
\begin{itemize}
  \item \emph{MVP (Woche 1\,--\,8):} Plattform-Templates fertigstellen; zwei Pilotteams onboarden; SLO-Standard-Set live; drei Services extrahieren; Canary-Pipeline verf\"ugbar.
  \item \emph{Scale (Woche 9\,--\,16):} weitere f\"unf Teams; w\"ochentliches Release; Error-Budget-Mechanik aktiv; Architekturboard mit ADR-Praxis.
  \item \emph{Messkonzept:} Outcomes (Lead Time, Change Failure Rate, MTTR, Audit-Findings) \textendash{} keine Aktivit\"atsmetriken (``wie viele Tickets bearbeitet'').
\end{itemize}

\smallskip
\textbf{6. Zwei Entscheidungen unter Unsicherheit (Pflicht):}
\begin{enumerate}
  \item \emph{Welche drei Services zuerst extrahieren?} \\
    Wahl: \emph{Claim Intake}, \emph{Fraud Check}, \emph{Notifications}. Begr\"undung: Hoher Lasthebel (Intake), klare Capability (Fraud), wenig Datenkopplung (Notifications). \emph{Verworfene Alternative:} \emph{Payout} zuerst \textendash{} zu hohes Audit-/Finance-Risiko, kein guter Lern-Service. \emph{Fr\"uhwarnsignale:} Wenn Lead-Time f\"ur \emph{Intake} sich nach 4 Wochen verdoppelt oder Fraud-Score-Genauigkeit sinkt \textrightarrow{} Schnitt \"uberpr\"ufen.
  \item \emph{Welche Release-Policy f\"ur Kernprozesse?} \\
    Wahl: \emph{Zweiw\"ochentlich + Canary 5\,\% / 25\,\% / 100\,\%}, Rollback automatisch bei Fehlerquote $>$\,1\,\%. Begr\"undung: Kernprozesse $=$ auditrelevant + finanzkritisch; w\"ochentlich nicht durchhaltbar bei aktuellem Reifegrad, monatlich w\"are zu langsam. \emph{Guardrails:} kein Release am Freitag; Security-Approval Pflicht; Rollback-\"Ubung monatlich. \emph{Fr\"uhwarnsignal:} Wenn Change Failure Rate im Kern $>$\,15\,\% in 2 Iterationen \textrightarrow{} Frequenz zur\"uck auf monatlich.
\end{enumerate}

\smallskip
\textbf{Anschluss an Strategie:} Auditierbarkeit (Compliance), Time-to-Market (Peak-Saison), Skalierbarkeit (mehrere Standorte). Die Plattform ist kein IT-Projekt, sondern ein \emph{steuerbares System}, das Entscheidungen, Risiko und Verantwortung sichtbar macht.
\end{loesungbox}

\clearpage

% --- Abschluss -----------------------------------------------------------
\section*{Abschluss}
\thispagestyle{plain}

\noindent Die Musterl\"osungen sind eine Diskussionsbasis. Wenn Ihre eigene L\"osung anders ist, fragen Sie: \emph{Habe ich andere Annahmen \textendash{} oder einen anderen Service-Schnitt angelegt?} Beides ist legitim, solange die Logik nachvollziehbar bleibt.

\medskip
\begin{infobox}[N\"achster Schritt]
Bringen Sie Ihre eigenen L\"osungen in die Coaching-Sitzung mit. Im Fokus stehen drei Fragen: Welche Annahmen haben Sie gemacht? Welche Zielkonflikte (Speed vs.\ Stabilit\"at, Autonomie vs.\ Standardisierung, Aufwand vs.\ Auditierbarkeit) haben Sie sichtbar gemacht? Welche Entscheidung haben Sie getroffen \textendash{} und w\"urden Sie sie auch verantworten, wenn die n\"achste Reorganisation kommt?
\end{infobox}

\vspace{1cm}
\noindent{\color{lpAccent}\rule{\linewidth}{0.6pt}}\\[4pt]
{\footnotesize\sffamily\color{lpGray}Lernplatt \textperiodcentered{} by Can Siebert \textperiodcentered{} Kurs 71 (Dig 42) \textperiodcentered{} Tag 10 \textperiodcentered{} L\"osungsheft \textperiodcentered{} \textcopyright{} 2026}

\end{document}
